\documentclass[12pt]{article}
\usepackage[utf8]{inputenc}
\usepackage[T1]{fontenc}
\usepackage{hyperref}
\title{Usoi Language Resources for K-12 Curriculum Development}
\author{Compiled by Claude Sonnet 5 (claude-sonnet-5) for Andr\'as Kornai}
\date{}
\begin{document}
\maketitle
\section{Introduction}
Usoi (Usui, Uchoi; ISO 639-3 usi) is a southern variety of Tripura (Kokborok), a Tibeto-Burman language of the Boro-Garo group, spoken in the southern Chittagong Hill Tracts, Bangladesh. Speaker figures differ between sources: about 25,000 (PARADISEC) versus about 9,400 for the Uchoi Tripura people group (Joshua Project).
\section{Primary Educational Resources}
SIL Bangladesh reports a mother-tongue literacy programme for Usoi adults and a multilingual education programme for children in development. No actual textbook was found.
\begin{itemize}
\item SIL Bangladesh document: \url{https://www.silbangladesh.org/sites/ban/files/the_light_which_can_illuminate_the_future.pdf}
\end{itemize}
\section{Linguistic and Reference Materials}
\begin{itemize}
\item Huziwara, phonology of the Usoi Tripura dialect (Journal of Asian and African Studies): \url{https://tufs.repo.nii.ac.jp/record/2054/files/jaas077005.pdf}
\item Joshua Project language page: \url{https://www.joshuaproject.net/index.php/languages/usi}
\end{itemize}
\section{Religious and Community Resources}
None found in search results.
\section{Digital and Audio Resources}
\begin{itemize}
\item PARADISEC collection MMR1: \url{https://catalog.paradisec.org.au/repository/MMR1}
\item Tripura recordings listing at Research Data Australia: \url{https://researchdata.edu.au/tripura-recordings/1700838}
\end{itemize}
I only saw these in search snippets and did not open them, so I cannot confirm they contain Usoi material.
\section{Curriculum Development Context}
The language is not known to be taught in schools. Kokborok materials from India may be adaptable, but I did not search for them. The SIL Bangladesh programme is the most relevant lead for mother-tongue education.
\end{document}